% chapters/frontmatter.tex - 题献与原书前言

% ==================== 题献 (Dedication) ====================
\thispagestyle{empty}
\vspace*{7cm}
\begin{flushright}
    \itshape
    献给我的妻子与儿子们:\\
    Sara, Georg 与 Lukas
\end{flushright}
\vfill
\cleardoublepage

% ==================== 译者序 ====================
\chapter*{译者序}
\addcontentsline{toc}{chapter}{译者序}

微局部分析 (Microlocal Analysis) 起源于 20 世纪中叶偏微分方程, 拟微分算子与现代奇性分析的交汇点. 从 Sato (佐藤干夫) 对超函数与微局部性的构想, 到 Hörmander (赫尔曼德), Duistermaat, Melrose 等人建立的经典伪微分算子与 Fourier 积分算子理论, 这一学科将偏微分方程的研究舞台从传统的位形空间底流形 $M$, 彻底提升到了经典力学与辛几何的天然载体——相空间 (余切丛 $\TstarM$). 在相空间中, 偏微分算子的主符号对应着经典力学中的 Hamilton 能量函数, 而算子解的奇性传播, 则精确化为相空间中奇性 (波前集 $\WF(u)$) 沿着 Hamilton 向量场流线的演化.

Peter Hintz 教授所著的这本《微局部分析导论》(\emph{An Introduction to Microlocal Analysis}, Springer GTM 304, 2025), 是近年来该领域不可多得的杰作. 作者师从 András Vasy 教授, 在广义相对论时空几何, 黑洞稳定性分析以及双曲型方程的微局部散射理论中做出了大量奠基性贡献. 与传统的百科全书式专著不同, 本书具有以下鲜明特色:
\begin{enumerate}
    \item \textbf{自成体系, 步步为营}: 除基本泛函分析和初等微分流形外, 全书从 $\R^n$ 上的 Schwartz 函数与缓增分布起步, 完整, 细致地证明了伪微分算子代数, 流形上的坐标不变性, 波前集几何以及实主型奇性传播的所有关键定理.
    \item \textbf{引入现代前沿工具}: 本书系统阐述了径向集上的奇性传播 (Radial Point Estimates) 与微局部 Fredholm 理论——这是自 Melrose 和 Vasy 以来现代散射理论与双曲 PDE 求解的核心利器.
    \item \textbf{全景式的真实物理与几何应用}: 在全书最后一章 (第~\ref{ch:11}~章), 作者综合运用全书建立的全部微局部工具, 完整给出了 de Sitter 空间上线型波动方程迟时渐近展开 (Late-Time Asymptotics) 与谐振态展开的严格证明.
\end{enumerate}

本中文重排版采用李文威老师开源的《代数学方法》书籍模板体系 (\texttt{AJbook.cls}) 重排. 在翻译过程中, 我们严格遵循数学专业规范与用户指引, 全书忠实于原著文字与证明推导, 采用英文半角标点符号, 保证全书交叉引用, 公式编号, 定理环境与参考文献的严密闭环.

希望本译本能为有志于现代偏微分方程, 微局部分析, 广义相对论及微分几何的研究人员与研究生提供一份严谨清晰的中文研读参考.

\vspace{1.5em}
\hfill \textbf{译者}\quad 谨识于 2026 年

\cleardoublepage

% ==================== 原书前言 (Preface) ====================
\chapter*{前言}
\addcontentsline{toc}{chapter}{前言}

在当前的时代背景下, 微局部分析大体上已不再是一门仅为自身而存在的研究领域, 而更是一套体量庞大, 功能强劲, 模块化且多面适用的工具集, 伴随着分析线性偏微分方程 (PDE) 的极具穿透力的新视角. 它在广义相对论, 混沌动力系统, 谱渐近分析以及反问题等众多截然不同的领域中, 都取得了令人瞩目的深刻应用, 这里仅列举数例.

微局部分析的一个特别核心的落脚点, 在于研究偏微分方程解的奇性 (singularities). 与传统的局部分析 (关注分布在特定关注点附近的局部正则性) 不同, 微局部分析关注的是在相空间 (phase space, 即位置与动量空间) 中的正则性. 相空间亦是经典力学中 Hamiltonian 力学的展开舞台; 每一个线性偏微分算子, 都在相空间上诱导出一个 Hamiltonian (即``能量函数''). 微局部分析提供了一整套工具, 能够将关于底层经典力学系统的信息, 翻译转化为对偏微分方程的定量估计. 由 Hörmander 建立的关于奇性传播的奠基性定理 (将于第~\ref{ch:08}~章证明), 便是这一深刻交融的优美范例: 偏微分方程解的奇性, 精确沿着底层经典力学系统的相空间轨道传播. 微局部分析最核心的算子工具, 是伪微分算子 (pseudodifferential operators, 常简记为 ps.d.o.) 代数. 伪微分算子允许人们将分布在相空间中进行局部化. 同时, 它们作为椭圆型偏微分算子的逆算子或拟逆 (parametrix, 即近似逆), 在分析中自然且独立地产生. 本书前半部分的几乎全部篇幅, 都致力于在欧氏空间及流形上完整, 详尽地建立伪微分算子理论.

在第~\ref{ch:11}~章中, 我通过完整证明 de Sitter 时空上线型波动方程解在迟时 (late-time) 渐近展开这一``现实世界''场景, 展示了本书所建立的几乎全部微局部方法的巨大威力. 许多更经典以及更新近的论题未能在本书的篇幅内展开, 其中包括: 边界值问题, 在拟线性偏微分方程中的应用, 束缚集 (trapped sets) 附近的估计 (在广义相对论中尤为关键), 非紧与奇异空间上的微局部分析 (具界几何, 几何微局部分析, $b$-演算), 波在奇异空间 (边界, 角柄, 锥, 棱边) 上的奇性传播与衍射, Fourier 积分算子, 谱渐近分析, 半经典缺陷测度, 变阶空间, 解析微局部分析, 仿微分技术等.

本书微局部分析理论的展开基本上从零起步. 唯一的预备知识是对泛函分析的初步了解. 此外, 读者若具备光滑流形理论的先验背景将大有裨益, 不过我们在第~\ref{sec:5.2}~节中对主要定义与几何构造给出了快速而完备的复习. 尽管微局部分析的核心思想高度概念化且呈现优雅的模块化结构, 但在构筑其底层技术支撑时仍需付出扎实深入的努力. 在本书中, 自第~\ref{ch:02}~章起, 几乎每一个数学结果都给出了极为详尽的完整证明. 我希望这种处理方式能够最大程度地降低研读门槛. 与此同时, 我在全书中也始终致力于凸显核心思想与主干脉络.

每章末尾都配有极为充实的习题. 部分习题要求读者验证正文中提及的若干推导细节或补足省略的论证; 另一些习题引导读者具体计算定义与定理的生动实例; 还有一部分习题则更具挑战性, 指导读者在提示下证明一些文献中频繁使用但未收录于本书正文的趣味结论.

本书是在我于麻省理工学院 (MIT) 2019 年春季学期主讲的研究生课程 18.157 (并在 2020 年及 2021 年苏黎世联邦理工学院 ETH Zürich 再次讲授) 的讲义基础上, 经过彻底修订与大幅扩展而成的. 根据主讲教师的教学目标与侧重点不同, 本书相当可观的一部分内容完全可以在一学期的研究生初阶课程中讲授. 例如, 一个自然的教学路线是: 详尽讲授伪微分算子与波前集的基本理论, 给出奇性传播定理的证明并勾勒若干代表性应用, 同时跳过所有半经典方法的章节.

本书后半部分 (特别是第~\ref{ch:10}~章与第~\ref{ch:11}~章) 涵盖了更具进阶性质的前沿专题, 富有进取心的授课教师完全可以讲授其中的部分内容. 我期望它们不仅能够成为有志于深入探索微局部领域的研究生与科研人员的起跳板, 也能凭借此处提供的充实细节, 作为一部可靠耐读的案头参考书. 各章末尾的``注记'' (Notes) 主要提供了相关原始文献的引用以及不同视角的处理途径.

在微局部分析相关领域, 已有许多优秀的讲义与专著, 各自侧重点不同. 在经典著作中, Hörmander 的著名巨著~\cite{Hor03,Hor05,Hor07,Hor09} 所涵盖的广度与深度远超本书, 但其定位绝大部分属于百科全书式的案头权威参考, 而非研究生导引教材; 此外, 第~\ref{ch:10}~章中自 Melrose~\cite{Mel94} 和 Vasy~\cite{Vas13} 以来在现代散射理论中扮演愈发核心角色的径向点估计, 在其中并未涉及. Taylor 关于线性与非线性偏微分方程的卓越著作~\cite{Tay81,Tay11a,Tay11b,Tay11c} 涵盖了极其宏富的素材, 重度使用了伪微分技术, 但狭义相空间奇性分析意义下的微局部分析在其中仅占次要地位. 其他重要教材包括: Shubin~\cite{Shu87} (侧重谱渐近分析), Alinhac--Gérard~\cite{AG07} (在 $\R^n$ 上引入 ps.d.o. 并迅速借助不动点定理切入非线性 PDE 的应用), 以及 Grigis--Sjöstrand~\cite{GS94} (强调 Fourier 积分算子与辛几何).

近年来带有半经典微局部特色 (涉及参数 $h \to 0$, 通常代表逆频率) 的教材包括 Zworski 关于半经典分析的教材~\cite{Zwo12}, 以及 Dyatlov--Zworski 关于散射共振数学理论的专著~\cite{DZ19}. 在本书中, 半经典概念的展开仅限于作为其经典对等物的自然变体; 而在上述两本专著中, 半经典方法被推向极深, 旨在研究高能极限下的谱与散射渐近行为. 特别地, 本书对~\cite[附录 E]{DZ19} 中汇集的核心材料给出了详尽阐述, 且本书第~\ref{ch:11}~章与~\cite[\S 5.7]{DZ19} 密切相关. Melrose 在该领域的经典讲义~\cite{Mel07} 深刻影响了我的行文思路, 尤其是本书前半部分. 我亦从 Wunsch~\cite{Wun13} 与 Vasy~\cite{Vas18} 的讲义中汲取了丰沛灵感; 后者在极大一般性下建立了散射伪微分演算, 并应用于渐近欧氏空间上的散射理论, 其中亦深入讨论了奇性传播与径向点估计.

我很幸运在哥廷根大学本科求学期间, 便在 Ingo Witt 教授精彩绝伦的偏微分方程系列讲座中领略了微局部分析的风采, 该课程的许多部分深受 Hörmander 著作的启发. 而在我学术生涯中施加最深远影响的, 毫无疑问是我的博士生导师 András Vasy 教授. 凡是熟悉其 PDE 研究风格的读者, 都能在本书的内容组织与数学品味中清晰辨认出他的深刻印记. 我对他在长年累月里毫无保留分享的远见卓识, 数学哲学与洞察力, 致以最由衷的崇高敬意. 此外, 我亦从众多师长, 同行与合作者处获益匪浅, 包括 Semyon Dyatlov, Jeffrey Galkowski, Rafe Mazzeo, Richard Melrose, Gunther Uhlmann, Jared Wunsch 以及 Maciej Zworski.

我由衷感谢修读我各门课程的全体同学, 以及提出诸多宝贵建议与修正的众多同行, 他们包括: Yonah Borns-Weil, André Dabrowski, Jesse Gell-Redman, Vincent Hoffmann, Jure Kališnik, Cosmin Manea, Aaron Moser, Pieter-Bart Peters, Phillip Pflaum, Jossi Schütt, Franziskus Steinert, Thomas Stucker, Ethan Sussman 以及 Tobias Weich. 我还要特别致谢 Jared Wunsch: 首先, 他率先在西北大学开设的研究生课程中勇敢试用了本书讲义的初版; 他随后反馈的宝贵意见极大提升了全书的明晰度与精确性. 其次, 正是他的热忱鼓励, 促使我将最初的讲义最终凝练成书.

在撰写本书期间, 作者得到了 Clay Research Fellowship, Sloan Research Fellowship, 美国国家科学基金会基金 DMS-1955614 以及瑞士国家科学基金会 SNSF Consolidator Grant TMCG-2\_223167 的资助.

\vspace{1.5em}
\hfill \textbf{Peter Hintz} \quad 谨识于 State College, USA

\cleardoublepage

% ==================== 利益声明 (Declarations) ====================
\chapter*{利益声明}
\addcontentsline{toc}{chapter}{利益声明}

\noindent\textbf{利益冲突}\quad 作者声明, 不存在与本手稿内容相关的任何利益冲突.
